\section{产品经理的底色：大厂游戏与长坡厚雪}

前面导读给出了 AI Native 的大问题，这一章先回到人。张月光对 AI 产品的判断，不是从模型研究出发，而是从十多年产品经理经历、两次创业和大厂组织游戏中长出来的。

张月光的经历很杂：SaaS、支付宝、字节、第一次创业、阿里优酷、妙鸭、第二次创业。访谈里他把 35 岁到 45 岁设成一个新阶段，希望和同一群人做同一件事。这解释了他为什么离开阿里：不是因为单个产品成功或失败，而是想摆脱大厂里的不可控变化，重新组织人和事。

\subsection{大厂游戏的两种战役}

他把大厂打怪升级分成两种：外部战役和内部战役。外部战役是赢市场、赢用户、赢竞争对手；内部战役是赢身边人、赢评价体系、赢资源分配。大厂增长期，电梯向上，外部战役很多；当行业进入真空期，内部战役的吸引力会变强，人很容易被组织机制拉回内部竞争。

\begin{warningbox}{内部战役的隐性成本}
内部战役可以让人升职，但会逐渐改变产品经理的注意力：从“用户和外部世界发生了什么”转向“组织里谁怎么看我”。AI 提供的新机会，恰恰是重新打开外部游戏。
\end{warningbox}

\subsection{第一次创业的情感线索}

第一次创业中的元音产品，让他第一次直接感到“活生生的用户”。用户在群里讨论、付费、停服时写小作文，这和大厂后台看到新增 500 万用户的抽象数字很不同。这个经历后来延续到星眠：AI 乙女游戏不是凭空出现，而是来自他早期对 Live2D、陪伴、角色和女性用户的长期兴趣。

\begin{knowledgebox}{从用户数字到用户生命感}
大厂数据能训练规模感，但不一定训练产品的情感判断。小产品的用户反馈会让产品经理看到具体的人、具体的喜欢和具体的失落，这会影响后续对 AI 陪伴和 AI 朋友的判断。
\end{knowledgebox}

\subsection{本章小结}

张月光的底层动力，是从内部游戏回到外部游戏，从抽象指标回到具体用户，从短期项目切换到长期同伴和长期事业。这个背景决定了他后面对 One Way Door、AI 朋友和 Agent 的偏好。

